% Plane-Wave Initial Condition: Spectral-Clean Construction on Periodic Grids
% Uses macros from preamble.tex

\section{Plane-Wave Initial Conditions on Periodic Grids}
\label{sec:plane-wave-ic}

\subsection{The spectral-leakage problem}
\label{sec:pw-leakage}

The plane-wave IC \texttt{--ic plane-wave --ic-wavevector k0} constructs
\begin{equation}
  \psi(x_i, t=0) = A \cos(k_0 x_i), \qquad
  \dot{\psi}(x_i, t=0) = A |k_0| \sin(k_0 x_i),
  \label{eq:pw-ic-cont}
\end{equation}
evaluated at the $N$ grid points $x_i = i L / N$, $i = 0, \ldots, N-1$ over a
periodic domain $[0, L]$. The rightward-propagating convention
($v = +|k_0| \sin$) is fixed by the CLI; see
\texttt{tidal/cli/\_simulate.py:\_plane\_wave\_slots}.

The discrete Fourier modes available on this grid are
\begin{equation}
  k_n = \frac{2\pi n}{L}, \qquad n = 0, 1, \ldots, N/2,
  \label{eq:discrete-modes}
\end{equation}
with $n = N/2$ being the Nyquist mode. The sampled field
$\cos(k_0 x_i)$ is \emph{monochromatic in the discrete Fourier
representation only when $k_0$ coincides with some $k_n$} --- that is, when
$k_0 L / (2\pi)$ is an integer below Nyquist. For any other $k_0$, the
continuous function $\cos(k_0 x)$ on $[0, L]$ is not actually periodic,
and the discrete Fourier transform of the sampled values has non-zero
amplitude on every mode, falling off as
$\approx 1 / (k_n - k_0)$ (spectral leakage). The FFT is faithfully
encoding the implicit boundary discontinuity that arises from treating
a non-periodic continuous function as periodic.

\textbf{Concrete example (dark-photon-plasma campaign configuration):} with
$L = 50$, $N = 32$, the discrete modes are $k_n \approx 0.126 \cdot n$.
Requesting $k_0 = 2.0$ gives a sampled IC whose rfft magnitude has
\begin{center}
\begin{tabular}{rrl}
$k$ & $|\hat{h}(k)| / \max$ & comment \\ \hline
$2.011$ & $1.000$ & $n = 16$ (Nyquist, peak) \\
$1.885$ & $0.0235$ & $n = 15$ \\
$1.759$ & $0.0113$ & $n = 14$ \\
$1.634$ & $0.0076$ & \\
$\vdots$ & $\vdots$ & \\
$0$     & $0.0022$ & DC leakage \\
\end{tabular}
\end{center}
Every discrete mode is excited above the $10^{-3}$ level.

\subsection{Why this matters: tachyonic amplification of leakage}
\label{sec:pw-tachyon-amplification}

For theories whose linearized dispersion is everywhere stable ($\mathrm{Re}(\lambda) \le 0$
for every generalized eigenvalue at every $k$), spectral leakage is
benign --- each leaked mode oscillates at its own frequency, contributing
at worst a small decoration to the observable.

For theories where the coupled system has tachyonic eigenvalues at some
$k$-modes, however, leakage onto those modes is \emph{exponentially
amplified} over the simulation interval. The modal solver's
perturbation-validity guard
(\texttt{tidal/solver/modal.py::\_evolve\_per\_mode},
\texttt{divergence\_threshold = 100}) then correctly rejects the run as
leaving the linearized regime, even though the intended physical IC
would have been stable.

The PGT dark-photon-plasma model (\texttt{examples/data/dark\_photon\_plasma.json})
exhibits exactly this failure mode. With $\delta_m \neq 0$ the
$h_5 \leftrightarrow a_1$ Gertsenshtein block merges with the torsion-trace
sub-block via the $\delta_m F_{ab} F_{\mathrm{torsion}}^{ab}$ kinetic mixing.
The merged block inherits tachyonic eigenvalues from the redundant
TorsionCDT state components at low $k$ ($\mathrm{Re}(\lambda) \approx 0.95$
at $k = 0$, decaying to zero near $k \approx 1$). Spectral leakage of
a $k_0 = 2$ IC onto these low-$k$ tachyons grows as
$\mathrm{e}^{\mathrm{Re}(\lambda) \cdot t}$ and trips the divergence guard
at $t = 50$ with amplitude ratio $\approx 8 \times 10^{13}$.

Two parallel concerns:
\begin{itemize}
  \item The \emph{IC spectral leakage} is a discretization artefact on the
        periodic grid, not physics. It can be eliminated at the IC
        construction stage (this document).
  \item The \emph{redundant TorsionCDT state with tachyonic mass squared}
        is a real feature of the full-tensor derivation (the
        fundamental-vector FV formulation projects these directions
        out at the Lagrangian level and does not have them). Documented
        and filed as a follow-up architectural issue; the spectrally-clean
        IC sidesteps the symptom without requiring the redesign.
\end{itemize}

\subsection{Fix: snap \texttt{--ic-wavevector} to the nearest discrete mode}
\label{sec:pw-snap}

The plane-wave IC constructor snaps each component of the requested
wavevector to the nearest discrete Fourier mode on periodic axes,
clamped below Nyquist:
\begin{equation}
  k_{\mathrm{snap}} = \Delta k \cdot
    \mathrm{clamp}\!\left(\round\frac{k_{\mathrm{req}}}{\Delta k},
                          -\left(\tfrac{N}{2}-1\right),
                          +\left(\tfrac{N}{2}-1\right)\right),
  \qquad \Delta k = \frac{2\pi}{L}.
  \label{eq:snap-formula}
\end{equation}
The Nyquist mode $n = N/2$ is excluded because $\cos(k_{N/2} x_i) = (-1)^i$
carries no directional information for real signals
($\sin(k_{N/2} x_i) \equiv 0$), and aliases with $n = -N/2$ under the
rfft representation. When
$k_{\mathrm{snap}} \neq k_{\mathrm{req}}$ by more than $1\%$ (relative), a
\texttt{Note:} is printed so the user knows what was actually simulated.

For non-periodic axes the Fourier basis is not the natural one
(Chebyshev, sine/cosine series, etc.), so no snap is applied on those
axes.

The velocity term $v = A |k| \sin(k \cdot x)$ uses $k_{\mathrm{snap}}$
for $|k|$ so the right-mover phase relationship is preserved.

\subsubsection{Override: \texttt{--ic-no-snap}}

Users who want the legacy verbatim-$k$ behaviour (e.g. reproducing a
pre-snap simulation, or testing a theory where off-grid wavevectors
carry physical meaning) can pass \texttt{--ic-no-snap} to disable the
snap. The flag is available on \texttt{tidal simulate},
\texttt{tidal sweep}, and \texttt{tidal sample}.

\subsection{Physical equivalence check: FV vs TorsionCDT}
\label{sec:pw-fv-vs-cdt}

With the snap enabled, the fundamental-vector dark photon model
(\texttt{torsion\_dark\_photon\_fv.json}, 10 fields) and the
TorsionCDT-plasma model
(\texttt{dark\_photon\_plasma.json}, 18 fields) give consistent
$P_{\max}$ at the campaign parameter point
$(m_A^2 = 0.955,\ \delta_m = 0.01,\ \xi = 0.274,\ \alpha_3 = 0.123)$ with the
equivalence mapping $\alpha_3 = m_T^2 / 2$. Observed values:
\begin{center}
\begin{tabular}{lll}
Model & $P_{\max}$ & relative difference \\ \hline
FV        & $0.001738$ & --- \\
TorsionCDT & $0.001759$ & $+1.2\%$ \\
\end{tabular}
\end{center}
Agreement at the percent level with two structurally different state
spaces confirms that the observables extracted from the TorsionCDT
formulation (despite its redundant dimensional components) match the
ghost-free FV formulation, as required by physical equivalence.

\subsection{When not to rely on the snap}

Three regimes where \texttt{--ic-no-snap} or careful grid design is
preferable:
\begin{itemize}
  \item \textbf{Sweeps in $k$}: if the user is scanning wavevector
        values as a physics parameter, snapping quantizes the scan
        onto the available discrete modes. For a clean continuous
        sweep either choose $L$ and $N$ commensurate with the target
        $k$ range, or accept the quantization with an explicit note.
  \item \textbf{Non-periodic simulations}: the snap only fires on
        periodic axes. Dirichlet/Neumann/absorbing axes need a basis
        appropriate to their boundary conditions; off-grid $k$ is not
        resolved cleanly regardless.
  \item \textbf{Gaussian IC with a travelling-wave factor}: gaussian
        IC also uses \texttt{--ic-wavevector}. At the time of writing
        the snap is applied only for \texttt{plane-wave}. Gaussian IC
        snap is filed as a follow-up.
\end{itemize}

\subsection{Implementation details}

Functions in \texttt{tidal/cli/\_simulate.py}:
\begin{itemize}
  \item \texttt{\_snap\_wavevector\_to\_grid(kvec, grid\_info, bounds)}:
        returns the snapped wavevector and a list of
        $(\mathrm{dim}, k_{\mathrm{req}}, k_{\mathrm{snap}})$ entries
        for dimensions where the snap changed $k$ by more than $1\%$
        (the list drives the user-visible \texttt{Note}).
  \item \texttt{\_plane\_wave\_slots(args, \ldots)}: applies the snap
        unless \texttt{args.ic\_no\_snap} is set.
\end{itemize}

Tests in \texttt{tests/test\_cli\_parsing.py::TestBuildInitialY0}:
\begin{itemize}
  \item \texttt{test\_plane\_wave\_snap\_off\_grid\_k} --- verifies
        $k_0 = 2.0$ at $N = 32, L = 50$ snaps to $n = 15$ ($k \approx 1.885$).
  \item \texttt{test\_plane\_wave\_snap\_on\_grid\_k\_unchanged} ---
        exactly on-grid $k$ passes through without change.
  \item \texttt{test\_plane\_wave\_snap\_disabled\_legacy\_behaviour} ---
        \texttt{ic\_no\_snap=True} uses $k$ verbatim.
  \item \texttt{test\_plane\_wave\_snap\_skips\_non\_periodic\_axes} ---
        non-periodic axes pass through.
  \item \texttt{test\_plane\_wave\_snap\_clamps\_below\_nyquist} ---
        requests well beyond Nyquist are clamped to the highest
        representable mode.
\end{itemize}
